

AI can be woven into apps in various exciting ways, each offering unique functionalities:

Natural Language Processing (NLP): By understanding human language, NLP can power chatbots that respond to user 
queries in a conversational manner, making interactions more natural and relatable.

Machine Learning (ML): ML enables apps to learn from user behavior and adapt over time. Whether it’s recommending 
products or optimizing performance, ML adds a layer of intelligence that makes apps more user-centric.

Computer Vision (CV): From recognizing faces to identifying objects, CV opens up a world of image-based interactions. 
Apps can see and understand the visual world, offering features that are both innovative and interactive.


PyTorch Mobile: An extension of the popular PyTorch framework, PyTorch Mobile is tailored for mobile platforms. 
It offers the flexibility to train and deploy machine learning models on mobile devices, providing a balance between 
performance and functionality.

Firebase MLKit: A robust set of machine learning APIs, Firebase MLKit equips developers with several pre-trained models 
suitable for tasks like image recognition, object detection, and text classification. Its seamless integration with 
Flutter ensures that developers can add AI features with minimal hassle.

TensorFlow Lite: A mobile-optimized version of TensorFlow, TensorFlow Lite is designed specifically for mobile devices. 
It allows developers to train and deploy machine learning models directly on mobile, ensuring efficient performance and 
reduced latency.

#buildinpublic - every day share results about the done work

Я много читала и изучала работы современных экспертов в вопросах саморазвития и личного роста: Gabor Mate, Yung Pueblo, 
Yuval Noah Harari, Thich Nhat Hanh, J. Krishnamurti. 

ArchOps

Solution Design


https://blockbuster.thoughtleader.school/p/dialectical-thinking-how-to-develop?utm_source=substack&utm_medium=email
Introducing The Art & Science Of Perspectiving

More than a specific practice, Perspectiving is a category of practices that share a common pattern:

    Value different perspectives than your own

    Identify a new and valuable perspective

    Embody that perspective and then look through it

    Extract insights from that perspective

    Name the perspective if it doesn’t have a name yet (optional)

    Repeat steps 2-5 over and over

    Collect broader insights about perspectiving as you go through steps 2-5

    Notice how your emerging perspective is different and better than others

    Use your unique perspectives to develop rare and valuable ideas

To gain deeper insight into what perspectiving looks like, below are…
6 Types of Perspectiving

    Dialectic Perspective. As I wrote about in Dialectical Thinking: How To Develop World-Changing Ideas, According To Research, simply understanding the pros and cons of each side of a conflict can yield tremendous insight. Dialectical Thinking counteracts the mind’s tendency to overgeneralize our current polarized perspective and to delete the best perspectives from the ‘other’ side.

    Sensorial Perspective. We can go directly to the phenomenon we’re studying and probe it with multiple senses. We can zoom in and zoom out. We can look at it from multiple locations. We can speed it up or slow it down. We can bring more senses to bear.

    State Perspective. We come up with radically different ideas in different states even though we’re the same person. As thought leaders, we can proactively alter our mental state in order to have unique insights. Many geniuses throughout history have tapped into unique insights while asleep, in a hypnogogic state, or in another state.

    Person Perspective. We can model someone’s internal and external thought patterns for their genius zone abilities. Then, we can look at the world through their perspective or even have a conversation with them in our minds. Alternatively, we can collaborate with others who are diverse to triangulate perspectives.

    Modern Polymathy. We can learn the tools, values, knowledge, theories, and methodology of a discipline and then apply that to solve a problem in a new way.

    Physics Perspective. We can understand how objects/entities in a discipline interact. Then, we can model that in our heads.

Now that we have a basic model of Perspectiving, let’s explore how many of history’s most creative people (Leonardo da Vinci, Thomas Edison, investing legend Ray Dalio, psychologist Edward de Bono, Einstein, and Picasso) used it as one of their fundamental creativity techniques…
